% Halle (Anatomical Institute Halle) human skull microCT -- pipeline narrative.
% Build:
%   pdflatex manuscript
%   pdflatex manuscript

\documentclass{compactarticle}

\usepackage{amsmath,amssymb,amsfonts,bm}
\usepackage{graphicx}
\usepackage[utf8]{inputenc}
\usepackage[T1]{fontenc}
\usepackage{booktabs}
\usepackage{enumitem}
\usepackage{xcolor}
\usepackage{float}
\usepackage{ragged2e}
\usepackage{url}

\graphicspath{{figures/}}

\providecommand{\br}{\toprule}
\providecommand{\mr}{\midrule}

\begin{document}
\justifying
\fancyhead[R]{\small\itshape Halle human skull pipeline}

\articletype{Validation report}

\title{MNI152 to Halle dry-skull microCT pipeline for human transcranial
focused ultrasound: cross-modality registration, ray-cast outer-table
extraction, and perpendicular-to-skull placement}

\author{Gianmarco F Pinton}

\affil{Joint Department of Biomedical Engineering, University of North
       Carolina at Chapel Hill and North Carolina State University}

\email{gia@email.unc.edu}

\keywords{transcranial focused ultrasound, human, micro-CT, skull, MNI152,
          ANTs SyN, angular spectrum, neuromodulation}

\begin{abstract}
We document the registration, surface-extraction, and placement steps that
together place a focused-ultrasound bowl on an MNI-defined anatomical target
through the Halle dry-skull human micro-CT (0.125\,mm isotropic, Zenodo
release). Three design choices distinguish this pipeline from its mouse
(DigiMouse $\rightarrow$ Maga 4K) and macaque (NMT v2) twins:
(i) cross-modality ANTs~SyN with a clinical-CT fixed image and a T1 MRI
template moving image (mutual-information cost), in place of the binary
cavity-to-brain-mask registration used for the mouse and macaque; (ii) a
ray-cast outer-bone-table surface that records the outermost cortical voxel
along each spherical ray from the bone-mass centroid, in place of a
morphological-closing envelope that would sit on the inner table;
(iii) perpendicular-to-skull bowl placement guided by surface normals and
constrained to a neighbourhood of a warped EEG site, in place of a
dorsal-cylinder apex search around the target's lateral projection. We also
document the empirically determined LAS-vs-LPS storage convention of the
Halle NRRD (the file header is wrong), the Hounsfield-anchored piecewise-
linear intensity-to-acoustic ramp used by the slab loader, and the
coordinate-frame migration required to bring this pipeline under TUBA's
``world-mm everywhere'' convention. Source data is reproducible: the Halle
volume is fetched from Zenodo by \texttt{setup/fetch\_halle\_skull.py}; the
MNI152 ICBM 2009a non-linear symmetric template is fetched by
\texttt{templates/fetch\_mni152.py}; processing modules live alongside the
neuromod\_parameters pipeline in \texttt{registration/}.
\end{abstract}

\section{Source data}\label{sec:source}
The Halle skull is a dry adult human cranium digitised by the Anatomical
Institute of the Martin-Luther-University Halle-Wittenberg and released on
Zenodo under \texttt{halle\_skull.nrrd} (1.6\,GB compressed)~\cite{halle_zenodo}.
The volume is reconstructed at a $0.125$\,mm isotropic voxel pitch from a
clinical-CT acquisition; intensities are calibrated in Hounsfield units
(HU), unlike the raw NRecon counts that the Maga 4K mouse and AMNH macaque
scans deliver. Air sits at $\mathrm{HU}\approx -1000$, soft tissue (when
present) at $\mathrm{HU}\approx 0$--$80$, and cortical bone at
$\mathrm{HU}\approx 700$--$2000$.

The reconstructed volume is dry skull only: no scalp, no in-situ brain, no
soft tissue. Downstream MNI-template registration therefore must rely on
the {\em outer} envelope of the bone (the head surface the MNI152
average-T1 volume sits on); the {\em inner} cavity, which the mouse and
macaque pipelines use, cannot be cleanly matched against MNI because
MNI152 is a head template (skull plus scalp), not a brain-only template.

Fig.~\ref{fig:halle_native} shows three orthogonal sections through the
\emph{native} 0.125\,mm Halle CT at the cortical-bone-mass centroid;
diplo\"e, trabecular basicranium, zygomatic processes, and dental
arches are resolved at the sub-millimetre scale that the microCT
acquisition was acquired for. The cavity is empty (air,
$\mathrm{HU}\approx -1000$). Table~\ref{tab:source} summarises the
scan parameters next to the DigiMouse, Maga 4K, and AMNH macaque
references for context.

\begin{figure}[H]
\centering
\includegraphics[width=\textwidth]{fig1_halle_native.png}
\caption{Halle dry-skull microCT at \emph{native} 0.125\,mm.
Orthoslices through the cortical-bone centroid (sagittal, coronal,
axial; left to right) at percentile-clipped $[0.5\%, 99.5\%]$
greyscale. Each panel is a $\sim$1\,150--1\,700-pixel cross-section
along its in-plane axes; cortical bone, diplo\"e, trabecular
basicranium, dental arches, and zygomatic processes are all resolved
at the sub-millimetre detail the microCT acquisition was designed for.
The endocranial cavity is air-filled (dry skull); the 1\,mm ANTs
working cache of \S\ref{sec:downsample} would smooth this detail down
to a $\sim$143-pixel grid.}
\label{fig:halle_native}
\end{figure}

\begin{table}[H]
\centering
\caption{Halle dry-skull microCT scan parameters in cross-species
context.}\label{tab:source}
\begin{tabular}{lllll}
\toprule
Property & DigiMouse & Maga 4K & AMNH macaque & Halle human \\
\midrule
Voxel pitch ($\mu$m, iso) & 100 & 6.127 & 60.6 & 125 \\
Field of view (mm$^3$, approx) & $38\times99\times21$ & $15.9\times23.8\times11.7$ & $130\times130\times130$ & $250\times250\times250$ \\
Coverage & whole body & head only (dry) & head only (dry) & head only (dry) \\
Soft tissue / in-situ brain & yes (cryo) & no & no & no \\
Intensity calibration & none (uint8) & none (uint16) & none (uint16) & HU (clinical CT) \\
Native file size & 78\,MB & $\sim$38\,GB & $\sim$5\,GB & $\sim$1.6\,GB \\
Atlas mode in TUBA & cavity\_binary & cavity\_binary & cavity\_binary & intensity \\
\bottomrule
\end{tabular}
\end{table}

The HU calibration is the practical entry-point for the human pipeline that
the mouse case did not have: it anchors the cortical-bone endpoint of the
piecewise-linear acoustic-property ramp (\S\ref{sec:acoustic}) to physical
density without any hydroxyapatite-phantom intervention, and it lets the
intensity-mode ANTs SyN cross-modality cost (\S\ref{sec:registration}) work
without first binarising the volume.

\section{MNI152 ICBM 2009a non-linear symmetric template}\label{sec:atlas}
The reference anatomy is the MNI152 ICBM 2009a non-linear symmetric
T1-weighted template~\cite{fonov2009}, a $1$\,mm isotropic average of 152
adult subjects produced by the McConnell Brain Imaging Centre. The
template ships with a companion brain mask
(\texttt{mni\_icbm152\_t1\_tal\_nlin\_sym\_09a\_mask.nii}) and is cached
in \texttt{templates/mni152\_icbm\_2009a/}; both are fetched by
\texttt{templates/fetch\_mni152.py}.

Two qualitative differences from the mouse Allen CCFv3 reference shape the
rest of the pipeline:

\begin{itemize}
\item {\em MNI is a head template.} The T1 voxels include scalp,
  subcutaneous fat, skull, and brain. There is no clean brain-only
  template against which a Halle cavity could be SyN-registered; the
  brain mask is provided as a separate post-segmentation product, not
  as the primary registration target.
\item {\em MNI is in RAS world coords.} Anatomical targets and EEG sites
  are conventionally quoted in MNI RAS millimetres. ANTs internally
  operates in LPS world coords, so the placement helper carries an
  explicit RAS$\leftrightarrow$LPS sign convention at each transform
  application step (\S\ref{sec:placement}).
\end{itemize}

We publish two named-coordinate dictionaries downstream alongside the
template binding:

\begin{itemize}
\item \texttt{ANATOMICAL\_TARGETS\_MNI} (50+ entries): standard atlas
  centroids in MNI RAS\,mm. Coverage includes S1, M1, V1, V5,
  the subthalamic nucleus, the ventral-intermediate (Vim) thalamus, the
  thalamus central, caudate, dorsolateral prefrontal cortex (DLPFC), ACC,
  PCC, insula, amygdala, entorhinal cortex, pre-SMA, the cerebellar
  vermis, auditory cortex (A1), temporal-lobe centroids, hippocampus,
  rPFC / lPFC / vmPFC, nucleus accumbens, and (as aliases) V1, A1, ACC,
  PCC, pre-SMA, cerebellum, and thalamus left/right pairs used by the
  human-skull sweep.
\item \texttt{EEG\_SITES\_MNI} (22 entries): standard 10--20 EEG site
  coordinates in MNI RAS\,mm, from Jurcak~et~al.~\cite{jurcak2007}.
  Fp1/2, F3/4/7/8, C3/4, Cz, CP3/4, P3/4, O1/2, Oz, T3/4/5/6, AF3/4.
\end{itemize}

These are consumed by the placement helper of \S\ref{sec:placement} to
resolve a (scalp-site, anatomical-target) pair into a bowl placement on
the Halle skull. The MNI152 binding lives at
\texttt{tuba.atlases.mni152.MNI152} with both dicts exposed at module level.

\paragraph{Anatomical annotation.}
For region-of-interest extraction, the MNI152 template is paired with
one of several volumetric parcellation atlases. The TUBA human binding
ships {\em three} atlases as drop-in alternatives, spanning the
coarse-to-fine spectrum that downstream analyses typically need:

\begin{itemize}
\item {\em Harvard-Oxford} (cortical-lateralized + subcortical, merged
  into 117 labels). The 1\,mm maximum-probability ``thr25'' variants
  distributed by FSL via the nilearn data cache~\cite{harvardoxford}.
  Coverage: 48 cortical regions $\times$ 2 hemispheres = 96 cortical
  IDs (Frontal Pole, Insular Cortex, ..., Occipital Pole), plus 21
  subcortical IDs (brainstem, plus left/right cerebral white matter,
  cerebral cortex, lateral ventricle, thalamus, caudate, putamen,
  pallidum, hippocampus, amygdala, accumbens). We merge the two with
  the subcortical IDs offset by the cortical maximum (96) to keep all
  117 IDs unique; the cortical-lateralized label is preferred where
  the two overlap.

\item {\em Schaefer 400 (7 networks)}. The 400-region cortical
  parcellation of Schaefer et al.~\cite{schaefer2018} aligned to the
  Yeo-7 functional networks, in MNI152 1\,mm space. 200 regions per
  hemisphere; subcortical structures are not covered.

\item {\em Schaefer 1000 (7 networks)}. The 1\,000-region variant of
  the same construction; the finest cortical parcellation we
  routinely use.
\end{itemize}

All three atlases live in MNI space on the FSL 1\,mm grid (shape
$182\times 218\times 182$, origin $(90, -126, -72)$ in RAS) which is
{\em different} from the ICBM 2009a template grid ($197\times 233
\times 189$, origin $(-98, -134, -72)$). The two grids are sampled at
the same world coordinates when figures need to overlay them.

The Glasser HCP-MMP1 cortical parcellation (360 cortical regions, 180
per hemisphere) is the natural high-resolution non-network alternative
but is distributed only via the BALSA database under registered access;
we omit it from the default binding for licensing reasons. The
swap-in pattern documented here (resample to MNI, then apply the
saved SyN forward transform with nearest-neighbour interpolation)
works for any volumetric label atlas in MNI space, so callers with
BALSA access can drop Glasser into the same slot.

These are the human counterparts of the Allen CCFv3 annotation in the
mouse pipeline: each provides a parcellation that is unambiguously
coordinate-aligned to the registration template and can be warped
through the saved SyN transforms to land in Halle skull space for
target-centroid extraction (\S\ref{sec:placement}).

Fig.~\ref{fig:mni_template} shows three orthogonal sections through the
MNI152 T1 template at the anterior commissure (world $(0,0,0)$) next
to the three parcellations at the same world coordinates; the layout
mirrors the mouse manuscript's template-and-atlas figure. The
granularity progression Harvard-Oxford~$\rightarrow$~Schaefer~400
$\rightarrow$~Schaefer~1000 is visible immediately: Harvard-Oxford
identifies anatomical lobes and major subcortical nuclei, Schaefer 400
breaks each cortical gyrus into a handful of patches aligned to
functional connectivity boundaries, Schaefer 1000 resolves the cortical
mantle at the network-cluster scale.

\begin{figure}[H]
\centering
\includegraphics[width=\textwidth]{fig2_mni_template.png}
\caption{MNI152 ICBM 2009a T1 template ({\em column 1}; greyscale,
1\,mm isotropic, brain mask in cyan) next to the three parcellation
atlases ({\em columns 2--4}): Harvard-Oxford (117 cortical +
subcortical labels), Schaefer 400 (cortical only, 7-network),
Schaefer 1000 (cortical only, 7-network). Labels are rendered in
distinct saturated colours from the \texttt{gist\_ncar} palette with
a deterministic ID$\rightarrow$hue shuffle so neighbouring IDs are
well separated. Orthoslices at the anterior commissure (world
$(0,0,0)$\,mm). All three atlases live on the FSL MNI152 1\,mm grid,
which is anatomically aligned to but does not share its voxel
indexing with the ICBM 2009a template grid; each is sliced at its own
voxel index for world $(0,0,0)$.}
\label{fig:mni_template}
\end{figure}

\paragraph{Three additional bundled parcellations.}
The TUBA human binding also ships three parcellations covering scopes
that the Harvard-Oxford + Schaefer pair does not: classical
whole-brain ROIs (AAL3), high-resolution deep-brain subcortex
(Pauli 2017), and coarse functional networks (Yeo 2011). All three
fetch via the same nilearn pattern except AAL3, whose upstream host
(\texttt{www.gin.cnrs.fr}) carries a TLS-cert misconfiguration; a
one-shot ``\texttt{curl -Lk}'' followed by \texttt{tar xzf} into the
nilearn cache directory is documented in
:mod:\texttt{tuba.atlases.mni152}. Fig.~\ref{fig:mni_extras} shows
the three at sagittal/coronal/axial sections through MNI AC.

\begin{itemize}
\item {\em AAL3 (Rolls 2020, extending Tzourio-Mazoyer 2002):} 166
  whole-brain labels covering cortex, subcortex, and cerebellum. The
  classical ROI atlas in fMRI work since the early 2000s;
  $\sim$15\,k citations make it the most-recognised parcellation in
  practice even where finer Schaefer or Glasser variants exist.
\item {\em Pauli 2017 deep-brain subcortex (Pauli et al., Sci Data
  2018):} 16 high-resolution subcortical regions at 0.7\,mm pitch
  (Pu, Ca, NAC, EXA, GPe, GPi, SNc, RN, SNr, PBP, VTA, VeP, HN, HTH,
  MN, STH). The thalamic and basal-ganglia detail that
  Harvard-Oxford lumps together as ``thalamus'' or ``putamen''
  blocks. Useful for deep-brain TUS targets (Vim thalamus, STN,
  nucleus accumbens, ventral pallidum) where coarser parcellations
  give a single big blob.
\item {\em Yeo 2011 (Yeo et al., J Neurophysiol 2011):} 7 cortical
  functional networks (visual, somatomotor, dorsal attention, ventral
  attention, limbic, frontoparietal, default mode). Used when the
  target is described as a network rather than a specific gyrus.
\end{itemize}

\begin{figure}[H]
\centering
\includegraphics[width=\textwidth]{fig5_extra_parcellations_mni.png}
\caption{Three additional bundled MNI-space parcellations:
{\em column 1} AAL3 (166 cortex + subcortex + cerebellum labels);
{\em column 2} Pauli 2017 (16 deep-brain subcortical labels at 0.7\,mm
native pitch, mostly empty outside the basal-ganglia / thalamic
band -- this is by design, the atlas only labels deep structures);
{\em column 3} Yeo 2011 7-network cortical parcellation. Orthoslices
at the anterior commissure (world $(0,0,0)$\,mm); same colour-shuffle
convention as Fig.~\ref{fig:mni_template}. The atlases are catalogued
in \texttt{tuba.atlases.mni152.PARCELLATIONS} alongside Harvard-Oxford
and Schaefer; \texttt{warp\_atlases.py} iterates over the catalogue to
produce a warped Halle-space companion for each (see
Fig.~\ref{fig:warped_extras}).}
\label{fig:mni_extras}
\end{figure}

\section{Orientation determination: the LAS-vs-LPS gotcha}\label{sec:orientation}
The Halle NRRD header tags the volume as
\texttt{space=left-posterior-superior} (LPS), the standard radiological
convention. The actual voxel storage, however, is LAS:

\begin{table}[H]
\centering
\caption{Halle NRRD voxel axes vs.\ their anatomical
direction.}\label{tab:halle_axes}
\begin{tabular}{lll}
\toprule
Halle voxel axis & Anatomical direction & Notes \\
\midrule
$+i$ (axis 0, increasing) & patient LEFT     & matches LPS header \\
$+j$ (axis 1, increasing) & patient ANTERIOR & {\em conflicts} with LPS header (header says posterior) \\
$+k$ (axis 2, increasing) & patient SUPERIOR & matches LPS header \\
\bottomrule
\end{tabular}
\end{table}

The conflict on axis~1 was verified by direct slice-viewer inspection of
the Halle NRRD on 2026-04-28: the slice that the header claims is the
most anterior in fact contains the foramen magnum and the occipital
ridge; the slice the header claims is most posterior contains the
nasal cavity and the upper jaw. The cause is unknown (a quirk of the
upstream scanner export or NRRD writer), but the consequence is
documented and the data is trusted over the header.

Encoded canonically: the affine that maps Halle voxel indices to
anatomically-correct RAS world\,mm is

\begin{equation}
\mathrm{halle\_affine\_to\_ras}(v) \;=\; \mathrm{diag}(-v,\; +v,\; +v,\; 1)
\label{eq:halle_affine}
\end{equation}

where $v$ is the voxel pitch in mm. Using $\mathrm{diag}(-v,-v,+v,1)$
(``proper LPS'' from the header) would put patient~anterior at NIfTI
$-y$ (RAS posterior), mirror-flipping the registration and breaking
every downstream MNI overlay. The encoded affine and the documentation
of the empirical disambiguation both live in
\texttt{registration/halle\_constants.py}; downstream code that writes
or reads the Halle volume {\em must} go through this module rather than
trusting the NRRD header field.

We do not run a per-slice bone-mass orientation probe analogous to the
mouse's. The clinical CT comes off the scanner with the patient already
in approximate anatomical orientation, so the only ambiguity to resolve
is the one axis-1 sign described above. The skull's gross bilateral
symmetry plus the unambiguous superior-inferior bone-density gradient
(cranial vault thin, basicranium thick) suffice to lock down the other
two axes without computation.

\section{Downsampling to the 1\,mm ANTs cache}\label{sec:downsample}
The native $0.125$\,mm Halle volume is moderate at $\sim$1.6\,GB
($\sim$820 million voxels for a typical adult skull). It is too large
to feed to ANTs SyN directly --- the deformation field would consume
$\sim$10\,GB at full pitch, and per-iteration cost is gradient-of-MI
dominated which scales with the active voxel count.

The pipeline therefore caches a 1\,mm Halle CT NIfTI
(\texttt{halle\_ct\_1mm.nii.gz}) at factor-8 block-max decimation, and
hands that to ANTs as the {\em fixed} image. Block-max (rather than
mean) preserves cortical-bone peak intensities, so the same
$\mathrm{HU}=700$ threshold used at native carries over to the cached
volume without recalibration.

The cache writes inherit the LAS$\rightarrow$RAS affine of
Eq.~\ref{eq:halle_affine} at the {\em new} voxel pitch (1\,mm). All
downstream visualisations of the warped MNI products (T1 template,
brain mask) sit on this 1\,mm grid; the slab loader of
\S\ref{sec:placement} samples the {\em native} $0.125$\,mm volume
directly, so the 1\,mm cache is used only by SyN and its QC.

We do not provide an intermediate 0.5\,mm cache analogous to the
mouse's 50\,$\mu$m intermediate. The 1\,mm ANTs cache + 0.125\,mm slab
sampler suffice for the existing case set and there is no use case for
an in-between resolution.

\section{No PCA pre-alignment}\label{sec:prealign}
The mouse Maga pipeline rotates its native volume by PCA-derived Euler
angles (manuscript \S5 of the mouse report; angles $(\alpha_x, \alpha_y,
\alpha_z) = (+8.07\,^\circ, +8.40\,^\circ, -19.51\,^\circ)$ for the
Maga 4K scan) so that all downstream caches share a clean axis-aligned
voxel grid centred at the cortical-bone centroid.

The Halle pipeline does {\em not} apply a PCA pre-alignment. The Halle
NRRD's voxel axes after the Eq.~\ref{eq:halle_affine} sign convention
{\em already} align with RAS to within a small (a few degrees) residual
rotation that is well within the capture range of ANTs SyN's
mutual-information stage. Forcing an explicit PCA pre-alignment would
introduce an order-1 \texttt{ndimage.affine\_transform} pass that
softens cortical-bone intensities by $\sim$25\,\% and would require
re-anchoring the Hounsfield thresholds of \S\ref{sec:acoustic}; the
softening is unavoidable for the mouse (whose skull sits at a large
angle relative to the index axes) but adds cost without benefit for the
Halle data.

The TUBA convention for the human is therefore that ``world-mm''
coordinates are in NIfTI RAS world {\em with the origin at NRRD voxel
$(0,0,0)$, not at the bone centroid}. This is a small relaxation of
the strict TUBA spec (mouse and macaque both centre on the bone
centroid after PCA); the bone centroid for Halle lives at world
$\approx (-130, +100, +100)$\,mm depending on the scan's per-subject
position. The species module documents this and provides
\texttt{nrrd\_voxel\_mm\_to\_halle\_ras} for legacy callers that still
carry NRRD-voxel-mm coordinates (the only difference is a sign flip on
the first component).

A legacy landmark-based orientation refinement
(\texttt{register\_orientation\_first.py}) exists, using a five-point
Umeyama fit on nasion, inion, the two pre-auriculars, and the vertex.
It is not invoked by the production pipeline because ANTs SyN handles
the same residual orientation drift directly, and the landmark
auto-detection produced noisier results than the SyN affine init in
side-by-side QC. The script remains in \texttt{registration/} for
reproducibility of historical figures but is not on the current code
path.

\section{Outer-bone-table extraction}\label{sec:surface}
Because MNI is a head template and not a brain-only template,
the human pipeline does not extract an endocranial cavity for
registration: the registration is intensity-vs-intensity (\S\ref{sec:registration}),
not mask-vs-mask. The surface mesh is needed for a {\em separate}
purpose: bowl placement. The transducer rim sits on the outer table
of the bone, the bowl's concave surface aims through the bone at the
target, and the apex sits a fixed distance outside the bone. To support
this placement geometry, the pipeline produces a dense outer-bone-table
point cloud with per-vertex normals.

Two extraction methods are available. They give noticeably different
surfaces, and only one of them is correct for transducer placement.

\subsection{Morphological close+fill ($\sim$22\,mm)}
\label{sec:surface_closefill}
\texttt{registration/extract\_halle\_surface.py} takes the
$\mathrm{HU}\geq 700$ bone mask at 1\,mm working pitch, applies a 22\,mm
morphological closing (a ball-shaped structuring element large enough
to bridge the foramen magnum ($\sim$30\,mm), the orbital openings, and
the choanae), fills any residual interior holes, and runs marching cubes
to produce a triangle mesh. This recipe is the standard
``inflate-the-skull-to-the-head-envelope'' trick used to manufacture a
fake scalp surface from a dry-skull CT for landmark-based registration
against an MNI head template.

The output mesh ({\em halle\_outer\_surface.npz}, $\sim$50k verts +
faces) is geometrically smooth and topologically simple, but its
vertices sit {\em inside} the cortical bone or on the inner table.
The 22\,mm closing intentionally inflates the bone mask to bridge the
large foramina; the resulting boundary is the convex-hull-ish envelope
of the bone, not the actual outer cortical surface. Using these
vertices for transducer placement would put the rim of the bowl 1--3\,mm
into the skull, which is not where a transducer can physically be
mounted.

The close+fill mesh is retained as the {\em fallback} input to
\texttt{halle\_placement.py} in case the ray-cast cache below has not
been generated, and as a registration-target proxy for the legacy
landmark-based scripts.

\subsection{Ray-cast outer table (the production path)}
\label{sec:surface_raycast}
\texttt{registration/extract\_outer\_bone\_surface.py} solves the
inner-table problem by ray-casting outward from the bone-mass centroid.
At a 1\,mm working pitch:

\begin{enumerate}
\item Compute the bone-mass centroid in storage-mm:
$(c_L, c_P, c_S) = \mathrm{mean}(\{\mathbf{i}\,v : \mathrm{HU}(\mathbf{i})\geq 700\})\cdot v$.
\item Sample the unit sphere uniformly via spherical coords:
$N_\mathrm{azim}=240$ azimuth steps in $[0, 2\pi)$, $N_\mathrm{elev}=120$
elevation steps in $[-\pi/2, +\pi/2]$, weighted by $\cos(\mathrm{elev})$
to suppress pole bunching. Total ray count: 28\,800.
\item For each direction $\mathbf{\hat d}$, walk outward from the
centroid in 1-voxel steps; record the position of the {\em outermost}
voxel along the ray whose intensity exceeds the bone threshold:
$\mathbf{p}_\mathrm{hit} = \mathrm{argmax}_{s} \{s \,:\,
\mathrm{HU}(\mathbf{c} + s\cdot v\cdot\mathbf{\hat d}) \geq 700\}$.
Rays that never re-enter bone after their last hit (i.e., that exit the
skull from the inside) are skipped.
\item Convert the hit cloud from storage-mm to TUBA Halle RAS via
$\mathbf{p}_\mathrm{ras} = \mathrm{diag}(-1,+1,+1)\cdot\mathbf{p}_\mathrm{hit}$.
\item Estimate per-vertex outward normals via local PCA: at each vertex,
find its $k=12$ nearest neighbours (cKDTree), fit a plane to the
neighbourhood, take the smallest-eigenvector direction, and orient the
sign so it points away from the centroid.
\end{enumerate}

The resulting NPZ ({\em halle\_outer\_bone\_surface.npz}, $\sim$24--26k
verts and matching normals) is the production placement target. It
captures the outer cortical surface to within $\pm 1$\,mm (the ray-step
size), it does not bridge across foramina (the bowl can be aimed across
the foramen magnum if needed without snapping its rim onto a virtual
mass), and the normals point physically outward.

A small fraction of rays ($\sim$1\,\%) record a hit at a stale cortical
fragment far from the main shell --- pieces of the basicranium that
detach from the cranial vault under the 1\,mm working-pitch
quantisation. These appear as outliers on the rendered mesh but do
{\em not} affect placement, because the placement helper restricts its
candidate set to an EEG-seeded sphere ($\S$\ref{sec:placement}) within
which the outliers do not sit.

In TUBA the raycast surface extractor lives at
\texttt{tuba.core.surface.raycast\_outer\_bone\_surface\_world\_mm};
it accepts an explicit storage-to-world rotation/sign matrix so the
ray-cast can be expressed in any species' aligned-native frame. The
human species module hard-codes Eq.~\ref{eq:halle_affine} for this
matrix.

\section{ANTs SyN intensity registration}\label{sec:registration}
The cross-modality registration matches the Halle 1\,mm CT volume
against the MNI152 1\,mm T1 template. The grids are the same physical
size ($\sim$1\,mm pitch, $\sim$200\,mm field of view) and largely
overlap once axis signs are correct (\S\ref{sec:orientation}); only the
intensity scales differ. The cost function therefore has to span the
cross-modality gap, which mutual information handles in a way that
sum-of-squares cannot.

\paragraph{Choice of fixed and moving.}
The Halle CT is held as the {\em fixed} image and the MNI T1 is the
{\em moving} image:

\begin{verbatim}
ants.registration(
    fixed=halle_ct_1mm,
    moving=mni_icbm152_t1,
    type_of_transform='SyN',
    reg_iterations=(40, 20, 10),
    aff_metric='mattes',         # Mattes mutual information
    syn_metric='mattes',
    syn_sampling=32,
)
\end{verbatim}

This is the {\em opposite} convention from the mouse Maga and macaque
NMT pipelines (which hold the atlas as fixed and the species's binary
cavity mask as moving, with a mean-squares metric). The reason is
practical, not principled: with the Halle CT as fixed, the warped MNI
products (T1, brain mask) land on the {\em CT's} 1\,mm grid, which is
the grid the slab loader uses for visual QC overlays; if MNI were
fixed and Halle were moving, the warped CT would land on the 1\,mm MNI
grid and a separate inverse pass would be needed for QC. The forward
and inverse transforms returned by ANTs cover both directions equally
either way, so the choice of which is ``fixed'' is a UX optimisation.

\paragraph{Saved artefacts.}
The forward + inverse transforms are persisted to
\texttt{registration/ants\_transforms/} and indexed by stable text
lists:

\begin{itemize}
\item \texttt{ants\_syn\_forward.txt}: applies image transforms in the
  direction MNI $\rightarrow$ Halle. Used by
  \texttt{ants.apply\_transforms} to bring MNI volumes onto the Halle
  CT grid (the MNI brain mask and T1 template are warped this way for
  QC).
\item \texttt{ants\_syn\_inverse.txt}: applies image transforms
  Halle $\rightarrow$ MNI. For point transforms (as opposed to image
  warps) ANTs applies the spatial inverse, so to map a {\em point} from
  MNI RAS into the Halle frame, one uses
  \texttt{apply\_transforms\_to\_points} with the {\em inverse}
  transform list. The placement helper takes this path.
\end{itemize}

\paragraph{Sanity probes.}
After SyN converges, six MNI anatomical landmarks
\big((0,0,0), $(0,86,-38)$, $(0,-113,0)$, $(0,0,90)$,
$(-85,-15,-44)$, $(85,-15,-44)$~mm in MNI RAS for AC, nasion, inion,
$\mathrm{Cz}$-proxy, left and right pre-auriculars\big) are warped into
Halle LPS and printed for visual sanity-check. After the
LAS-vs-LPS sign fix of \S\ref{sec:orientation}, AC lands at the centre
of the bony cranium, nasion lands at the nasal-bone tip, inion lands at
the external occipital protuberance, $\mathrm{Cz}$ lands at the cranial
vault apex, and the pre-auriculars land at the temporal-bone external
auditory canals. Any axis-flip remaining in the pipeline manifests as
a $\pm 50$--$150$\,mm offset on these probes; they pass for the current
release.

\paragraph{Warped brain volume.}
The MNI brain mask warps to a Halle brain region of approximately
1\,446\,mL. This is within the expected range for adult human cranial
cavities (1\,100--1\,500\,mL) and confirms the registration's
gross scale is correct; the residual local mismatch is below the SyN
deformation tolerance ($\sim$1--2\,mm) and is not separately
reported. Fig.~\ref{fig:skull_atlas} overlays the warped MNI T1
template (warm colour, masked to the warped MNI brain mask) on the
Halle CT (bone, greyscale); the cyan contour traces the warped brain
mask. The mask follows the inner table of the cortical bone in every
section, with no mass-effect leakage into the orbits, the nasal cavity,
or the foramen magnum.

\begin{figure}[H]
\centering
\includegraphics[width=\textwidth]{fig3_skull_atlas.png}
\caption{Post-SyN registration QC. Native 0.125\,mm Halle CT (bone,
greyscale) overlaid with the warped MNI T1 template (1\,mm,
masked to the warped MNI brain mask, inferno colour map) and the
warped brain mask (cyan contour). Orthoslices through the warped
brain centroid. The cortical-bone boundary is rendered at native
microCT pitch and the brain envelope (1\,mm, smoother) fills the
endocranial cavity to within the SyN deformation tolerance
($\sim$1--2\,mm); the cyan contour follows the inner table of the
cortical bone in every section without leaking into the orbits, the
nasal cavity, or the foramen magnum.}
\label{fig:skull_atlas}
\end{figure}

ANTs SyN at this resolution completes in $\sim$3--5\,min on a single
Threadripper core (vs.\ $\sim$70\,s for the mouse Maga binary
registration at 25\,$\mu$m, whose deformation field is much smaller).

\paragraph{Warped-atlas products.}
After SyN converges, we additionally warp the six MNI-space
parcellations catalogued in
\texttt{tuba.atlases.mni152.PARCELLATIONS} (the three primary atlases
of \S\ref{sec:atlas} plus the three extras of
Fig.~\ref{fig:mni_extras}) through the saved forward transform list
using nearest-neighbour interpolation (preserves integer label IDs).
Each parcellation lives on a slightly different MNI grid from the
ICBM 2009a template; ANTs resamples in physical world coordinates so
the source-grid mismatch is handled internally. The three primary
warped products land on the Halle 1\,mm grid as
\texttt{harvard\_oxford\_117\_in\_halle.nii.gz} (117 labels),
\texttt{schaefer\_400\_7nets\_in\_halle.nii.gz} (400 labels), and
\texttt{schaefer\_1000\_7nets\_in\_halle.nii.gz} (1\,000 labels). The
extras land as \texttt{aal3\_in\_halle.nii.gz} (166 labels),
\texttt{pauli\_2017\_in\_halle.nii.gz} (16), and
\texttt{yeo\_7nets\_in\_halle.nii.gz} (7); the catalogue-iterating
\texttt{warp\_atlases.py} produces all six in a single pass.
Fig.~\ref{fig:warped_atlases} overlays all three on the native
$0.125$\,mm Halle bone. The labels are coarse compared to the bone
(1\,mm warp grid vs.\ $0.125$\,mm bone) but the cortical envelope of
each parcellation conforms to the inner table of the cortical bone in
every section; no region leaks through the cortical shell into the
extracranial space or through the foramen magnum. The 117-vs-400-
vs-1\,000 granularity progression is visible at native bone scale:
Harvard-Oxford gives lobar-level patches, Schaefer 400 gives
gyrus-level patches that follow the actual Halle anatomy, Schaefer
1\,000 resolves cortical mini-clusters that span only a few mm.

\begin{figure}[H]
\centering
\includegraphics[width=\textwidth]{fig4_warped_atlases_on_halle.png}
\caption{Three MNI-space parcellations warped through the SyN
transform onto the native Halle CT (the mouse-manuscript Fig.~6
analogue). {\em Column 1}: bare Halle bone (greyscale, native
0.125\,mm) with the warped MNI brain mask drawn as a cyan contour.
{\em Columns 2--4}: same bone with the warped Harvard-Oxford,
Schaefer~400, and Schaefer~1\,000 labels overlaid at $\alpha=0.7$.
Orthoslices through the warped brain centroid. The label volumes are
at 1\,mm (the ANTs warp grid); the bone is at native 0.125\,mm so the
cortical-bone boundary stays sharp.
\\[2pt]
{\em On the apparent coverage gaps.} The labels do not cover the full
brain mask volume in any of columns 2--4, and the bone shows through
where they are absent. This is a parcellation-coverage property, not a
registration error. Numerically, 99.1\,\% of Harvard-Oxford labels sit
inside the warped MNI brain mask (0.9\,\% are edge artefacts; 13.3\,\%
of the brain mask is correctly unlabeled white matter and
brain-stem-boundary tissue that Harvard-Oxford simply does not
parcellate); 99.4\,\% of Schaefer~400 labels sit inside the brain mask
(the Schaefer construction is cortical-only, so the entire subcortex,
cerebellum, and brainstem are unlabelled by design). One additional
caveat: the FSL atlases live on the MNI152NLin6Asym grid while the SyN
transforms were learned from the ICBM~2009a template (\S\ref{sec:atlas}),
so warped labels carry a few-millimetre residual offset relative to the
ICBM-aligned brain mask --- below SyN deformation tolerance but visible
in the highest-resolution Schaefer~1\,000 boundaries. Use the
highest-resolution parcellation that the downstream analysis actually
requires: Harvard-Oxford for lobar-level region averages, Schaefer~400
for network-aligned gyrus patches, Schaefer~1\,000 when sub-gyral
detail matters.}
\label{fig:warped_atlases}
\end{figure}

\paragraph{Warped extras (AAL3, Pauli, Yeo) on the Halle skull.}
Fig.~\ref{fig:warped_extras} overlays each of the three extras on the
native 0.125\,mm Halle bone (same layout as
Fig.~\ref{fig:warped_atlases}). The Pauli column makes the deep-brain
detail visible directly on the native CT: the bilateral putamen and
caudate flank the lateral ventricles, the substantia nigra and
subthalamic nucleus sit ventral to the thalamus, and the nucleus
accumbens occupies its ventromedial position --- structures that the
Harvard-Oxford column of Fig.~\ref{fig:warped_atlases} aggregates into
a single ``Right Caudate / Right Putamen'' block. The Yeo column shows
the seven canonical networks as large cortical strips that fall
neatly inside the inner table of the bone.

\paragraph{AAL3 brain-envelope masking.}
AAL3 is the only one of the six bundled parcellations whose native
source is intrinsically L/R-asymmetric ($\sim$8\,\% more
right-hemisphere voxels than left, concentrated in the lateral
temporal lobe and right Crus 2). After warping into Halle the R-side
overhang spills past the warped MNI brain mask and into the cortical
bone shell: pre-fix, 1.7\,\% of AAL3 label voxels overlapped the
$\mathrm{HU}>700$ bone region and 3.7\,\% sat outside the brain mask
entirely, with 90\,\% of the overlap on the patient-R side. To bring
AAL3 in line with the other parcellations (all $\leq 0.2\,\%$ bone
overlap) we apply a post-warp brain-envelope masking step:
$\mathrm{aal3\_in\_halle}[\,\mathrm{brain\_mask} = 0\,] \leftarrow 0$.
This removes 42\,039 voxels (3.7\,\% of the pre-mask total) and brings
the bone-overlap fraction down to 0.19\,\% and the L/R asymmetry to
$-0.8\,\%$ (near-symmetric).

The masking is gated by a per-atlas flag
(\texttt{Parcellation.mask\_to\_brain\_envelope}) so it applies only
to AAL3 and is idempotent under re-running
\texttt{warp\_atlases.py}. Adding a future atlas with similar L/R
asymmetry (e.g., Glasser HCP-MMP1 if BALSA-fetched and resampled to
MNI) is a one-line catalogue addition; the warp script picks up the
flag without modification.

\begin{figure}[H]
\centering
\includegraphics[width=\textwidth]{fig6_extra_parcellations_on_halle.png}
\caption{Three additional warped parcellations on the native Halle CT
(same layout as Fig.~\ref{fig:warped_atlases}). {\em Column 1}: bare
0.125\,mm Halle bone with the warped MNI brain mask in cyan.
{\em Columns 2--4}: warped AAL3 (166 labels, whole-brain;
brain-envelope-masked --- see paragraph above), Pauli 2017
(16 deep-brain subcortical), Yeo 2011 (7 cortical functional
networks) overlaid at $\alpha=0.7$. The Pauli column gives the
deep-brain detail useful for Vim, STN, NAc, and pallidum TUS
targeting; the Yeo column gives the network-level framing useful for
default-mode / salience / dorsal-attention TUS protocols. Post-mask,
all three extras have $\leq 0.2\,\%$ label voxels overlapping bone
(matching Harvard-Oxford); Pauli is intrinsically at $0\,\%$ because
its labels stop well inside the brain mask. A few-mm offset from the
FSL-vs-ICBM grid mismatch remains visible at the highest resolution
for atlases on the FSL MNI152NLin6Asym grid.}
\label{fig:warped_extras}
\end{figure}

\section{Acoustic-property mapping}\label{sec:acoustic}
The Hounsfield calibration of the Halle volume permits an
{\em absolute}, deterministic piecewise-linear ramp from intensity to
acoustic properties $(c, \rho)$, with no reliance on histogram modes
the way the mouse pipeline must (which has no HU calibration). Both
endpoints are anchored on published cortical-bone literature:

\begin{table}[H]
\centering
\caption{Halle intensity-to-acoustic ramp (linear interpolation between
endpoints; clipped at $\mathrm{HU}\leq I_\mathrm{low}$ to water,
clipped at $\mathrm{HU}\geq I_\mathrm{high}$ to cortical
bone).}\label{tab:acoustic}
\begin{tabular}{lll}
\toprule
Anchor & Intensity (HU) & $(c, \rho)$ \\
\midrule
$I_\mathrm{low}$  & 700  & $(1540\,\mathrm{m\,s^{-1}}, 1000\,\mathrm{kg\,m^{-3}})$ \\
$I_\mathrm{high}$ & 1973 & $(2900\,\mathrm{m\,s^{-1}}, 2200\,\mathrm{kg\,m^{-3}})$ \\
\bottomrule
\end{tabular}
\end{table}

$I_\mathrm{low}=700$\,HU is the standard radiological lower bound of
cortical bone in soft-tissue-windowed clinical CT; cancellous bone
sits slightly below, marrow and trabeculae below that. Setting
$c$ and $\rho$ to water-equivalent here (rather than soft-tissue
intermediate values) is consistent with the Maga matched-fluid coupling
assumption: voxels below the bone threshold are treated as a continuum
of acoustic water, and any remaining soft-tissue contribution to the
focal field is captured by the downstream brain-attenuation block of
\texttt{validate\_transcranial.skull\_to\_phase\_screens}.

$I_\mathrm{high}=1973$\,HU is the upper endpoint of the Halle cortical-
bone distribution, calibrated empirically against the highest-intensity
voxel population in this scan. The literature endpoint at $(2900,
2200)$\,m\,s$^{-1}$ / kg\,m$^{-3}$ is bound to this scan's HU mode rather
than to a numerically arbitrary saturation value, mirroring the design
discipline of the mouse pipeline (\S8 of the mouse report); the
difference is that for the human we get the HU calibration ``for
free'' from the scanner, while the mouse had to choose the cortical
histogram mode by hand. The mouse uses
$\rho_\mathrm{bone}=1900$\,kg\,m$^{-3}$ (the conservative DigiMouse
literature value); the human uses $2200$\,kg\,m$^{-3}$, slightly higher
to match the HU-calibrated cortical density.

Voxels in $[I_\mathrm{low}, I_\mathrm{high}]$ interpolate linearly
between the two endpoints. The dry-skull cavity of the Halle release
contains air and is below $I_\mathrm{low}$; it is therefore mapped to
water without any explicit cavity-infill step. For live-subject
clinical CT (which the Halle pipeline could in principle ingest with
no code change beyond the source path), the cavity would contain brain
parenchyma at $\mathrm{HU}\approx 30$--$50$, still below $I_\mathrm{low}$,
so the same ramp maps it to water --- consistent with how acoustic
solvers usually treat brain parenchyma in transcranial simulations at
clinical frequencies (250\,kHz--1\,MHz).

The ramp is implemented in
\texttt{tuba.core.slab.IntensityToAcousticRamp}, bound by the human
species module as

\begin{verbatim}
SLAB_RAMP = IntensityToAcousticRamp(
    i_low=700.0, i_high=1973.0,
    c_water=1540.0, rho_water=1000.0,
    c_bone_max=2900.0, rho_bone_max=2200.0,
)
\end{verbatim}

so the slab loader of \S\ref{sec:placement} produces $(c\_map,
\rho\_map)$ at the slab's sampling pitch with no further setup.

\paragraph{Attenuation.}
We carry over the frequency-power-law model from the Maga pipeline
unchanged: $\alpha_\mathrm{bone} = 8$\,dB\,cm$^{-1}$\,MHz$^{-1}$,
$\alpha_\mathrm{brain} = \alpha_\mathrm{tissue} =
0.5$\,dB\,cm$^{-1}$\,MHz$^{-1}$, with linear frequency dependence
($\alpha \propto f$). At a typical clinical 0.5\,MHz the per-mm
cortical-bone attenuation is $\sim$0.4\,dB; for a 5--10\,mm bone
traverse this is 2--4\,dB and dominates the transcranial loss budget
together with the impedance-mismatch reflections at the bone
interfaces.

\section{Placement: perpendicular-to-skull bowl
geometry}\label{sec:placement}
The placement helper takes (i) an anatomical target name resolving to
an MNI RAS coordinate, (ii) optionally an EEG-site name resolving to a
second MNI RAS coordinate, (iii) a focal length $F$\,(mm), and (iv) a
bowl aperture radius $R$\,(mm); it returns the bowl apex position,
the inward beam direction, and a 3D orthonormal beam basis. The
output is in TUBA Halle RAS world\,mm.

\paragraph{Coordinate-chain.}
For each input MNI RAS coord (target, EEG site) the helper applies the
chain
\begin{align}
\mathrm{MNI\,RAS} &\xrightarrow{(-,-,+)} \mathrm{MNI\,LPS} \nonumber\\
&\xrightarrow{\texttt{ants.apply\_transforms\_to\_points(inv)}}
   \mathrm{Halle\,LPS} \nonumber\\
&\xrightarrow{(-,-,+)} \mathrm{Halle\,RAS\,world\,mm}\label{eq:coord_chain}
\end{align}
where the first and last steps negate the first two components and the
middle step applies ANTs' inverse transform (since for points the
spatial inverse of the image registration applies). This is the
canonical TUBA convention for the intensity-mode atlas binding:
{\em images} use the saved forward transform list to go atlas $\rightarrow$
subject; {\em points} use the inverse list to go atlas $\rightarrow$ subject.

\paragraph{Surface candidate set.}
The outer-bone-table point cloud and per-vertex normals
(\S\ref{sec:surface_raycast}) are loaded once into a cKDTree and cached.
If an EEG seed is supplied, candidate vertices are restricted to those
within \texttt{eeg\_search\_radius\_mm} of the warped seed
($10$\,mm by default; widens to $20$\,mm if fewer than $50$ candidates
are found inside the initial radius). Without an EEG seed, all $\sim$25k
vertices are candidates.

\paragraph{Perpendicular fit.}
For each candidate vertex $\mathbf{p}_i$ with outward normal $\mathbf{n}_i$,
the perpendicular distance from the inward-normal line
$\{\mathbf{p}_i + s(-\mathbf{n}_i)\,:\,s\in\mathbb{R}\}$ to the target
$\mathbf{T}$ is
\begin{equation}
d_i = \big\|\,(\mathbf{p}_i - \mathbf{T}) -
          \big[(\mathbf{p}_i - \mathbf{T}) \cdot (-\mathbf{n}_i)\big](-\mathbf{n}_i)\,\big\|
\end{equation}
We additionally require the target to be ``in front'' of the
transducer --- the dot product $(\mathbf{p}_i - \mathbf{T})\cdot(-\mathbf{n}_i)
< 0$ --- which excludes candidates on the opposite side of the head
that happen to have a normal nearly parallel to the target direction.
The selected vertex is the one minimising $d_i$ among in-front
candidates.

\paragraph{Bowl geometry.}
The selected vertex is the bowl rim contact. The bowl apex sits
$h$\,mm outside the bone along $-\mathrm{beam}$, where
\begin{equation}
h \;=\; F \;-\; \sqrt{F^2 - R^2}
\end{equation}
is the geometric depth of a spherical-cap bowl of focal length $F$ and
aperture radius $R$. For the default $(F, R) = (30, 15)$\,mm,
$h = 4.02$\,mm. The focal point sits at apex~$+ F\cdot\mathrm{beam}$,
typically $\sim$0--3\,mm from the warped MNI target (since the bowl rim
is constrained to the outer skull, not to a sphere of radius $F$ around
the target; the perpendicular-fit residual $d_\mathrm{best}$ is the
target-to-beam-axis miss).

\paragraph{Output basis.}
The placement helper returns a full 3D orthonormal frame:
$\mathbf{n}_3 = \mathbf{\hat b}$ (the inward beam direction),
$\mathbf{t}_3$ orthogonalised against $\mathbf{n}_3$ from the seed
$(1,0,0)$ if the beam is not nearly L-axis-aligned else from $(0,1,0)$,
$\mathbf{b}_3 = \mathbf{n}_3 \times \mathbf{t}_3$. A legacy 2D
projection $(n_\mathrm{LP}, t_\mathrm{LP})$ in the LP plane is also
included for back-compatibility with older slab consumers.

In TUBA the placement helper lives at
\texttt{tuba.core.placement.place\_perpendicular\_to\_surface}; the
human species module's \texttt{place\_on\_skull} wraps it with the MNI
target lookup, the EEG site lookup, the coordinate-chain of
Eq.~\ref{eq:coord_chain}, and back-compat keys carrying the
NRRD-voxel-mm coordinates for legacy callers.

\paragraph{Sample placement: CP3 $\rightarrow$ S1\_left.}
A representative case in default geometry: target \texttt{S1\_left} at
MNI RAS $(-40, -25, 55)$, EEG seed \texttt{CP3} at MNI RAS
$(-55.9, -44.5, 56.8)$, $F=30$\,mm, $R=15$\,mm,
\texttt{eeg\_search\_radius\_mm}~$=10$\,mm. The resolved Halle-RAS
quantities are listed in Table~\ref{tab:placement_sample}.

\begin{table}[H]
\centering
\caption{Resolved Halle placement for CP3 $\rightarrow$
S1\_left.}\label{tab:placement_sample}
\begin{tabular}{ll}
\toprule
Quantity & Value (Halle RAS world\,mm) \\
\midrule
S1\_left (warped target)        & $(-109.5,\, +82.3,\, +133.9)$ \\
CP3 (warped EEG seed)           & $(-145.5,\, +66.2,\, +138.8)$ \\
Bowl rim on bone                & $(-126.1,\, +72.9,\, +149.9)$ \\
Bowl apex (source plane)        & $(-128.8,\, +71.8,\, +152.7)$ \\
Inward beam direction           & $(+0.665,\, +0.276,\, -0.694)$ \\
Bowl depth $h$                  & $4.02$\,mm \\
Perpendicular residual          & $2.79$\,mm \\
Apex-to-target distance         & $29.0$\,mm \\
EEG-seed-to-placement offset    & $19.7$\,mm \\
\bottomrule
\end{tabular}
\end{table}

The apex-to-target distance ($29.0$\,mm) is within $1$\,mm of the
focal length ($30$\,mm), so the geometric focus and the anatomical
target coincide to within $\sim$$1$\,mm at this placement; the
perpendicular residual ($2.79$\,mm) is the off-axis miss of the beam
axis from the target. The EEG-seed-to-placement offset is large
($19.7$\,mm) because the CP3 location is on scalp and the bowl rim is
on bone, but the rim is on the skull region under CP3 (the lateral
parietal vault) within the EEG-constrained search radius.

\section{Slab loader and acoustic-property maps}\label{sec:slab}
Given a placement (apex, beam direction) and a slab size (lateral,
elevation, depth in metres) and pitch $\Delta x$, the slab loader
samples the {\em native} 0.125\,mm Halle NRRD on a beam-aligned grid:

\begin{itemize}
\item Slab origin is the apex (i.e., $z=0$ at the source plane).
\item Slab axes are lateral $(\mathbf{t}_1)$ and elevation $(\mathbf{t}_2)$
  orthogonal to the beam, built by Gram-Schmidt from a helper seed of
  $(0,0,1)$ unless the beam is near-vertical
  ($|\mathrm{beam}_z| > 0.9$), in which case $(1,0,0)$.
\item Sampling is via
  \texttt{scipy.interpolate.RegularGridInterpolator} on the Halle NRRD,
  with world-axis arrays built from Eq.~\ref{eq:halle_affine} so that
  axis 0 ranges monotonically downward from $0$ (axis 0 of the NRRD is
  patient-L, which TUBA-Halle-RAS sees as $-x$).
\item Intensities are passed through the ramp of
  Table~\ref{tab:acoustic} to produce $(c\_map, \rho\_map)$ at the slab
  pitch.
\item No cavity-infill step is applied (\S\ref{sec:acoustic}): the
  dry-skull cavity is at HU below $I_\mathrm{low}$ and the ramp maps it
  to water.
\end{itemize}

The slab loader and the placement helper both live in TUBA's
\texttt{tuba.species.human} module; downstream callers (the
angular-spectrum solver, the validation-report builder) consume the
returned $(c\_map, \rho\_map, dz\_slab, frame)$ tuple identical in
shape to the mouse and macaque loaders.

\section{TUBA migration and lessons carried from the mouse case}
\label{sec:tuba}
TUBA (Transcranial Ultrasound Brain Atlas) is a unifying library that
exposes the four-stage pipeline --- orientation $\rightarrow$ align
$\rightarrow$ surface/cavity $\rightarrow$ ANTs SyN --- as
species-agnostic primitives in \texttt{tuba.core}, with thin species
modules under \texttt{tuba.species} that bind numerics (voxel size,
thresholds, atlas) to one of three references (mouse Allen CCFv3,
macaque NMT v2, human MNI152). The human module is
\texttt{tuba.species.human}; its core dependencies are
\texttt{tuba.core.\{surface, warp, placement, slab\}} plus
\texttt{tuba.atlases.mni152}.

\subsection{Three design choices carried forward unchanged}
\begin{itemize}
\item {\em ANTs convention frozen.} For the mouse and macaque, the
  registration is
  \texttt{ants.registration(fixed=atlas\_mask, moving=cavity\_mask)};
  for the human it is
  \texttt{ants.registration(fixed=halle\_ct, moving=mni\_t1)}. In both
  cases image forward transforms bring the {\em moving} image onto the
  {\em fixed} grid, and image inverse transforms reverse the
  direction; point transforms via
  \texttt{ants.apply\_transforms\_to\_points} apply the spatial inverse
  of the corresponding image warp. Users of TUBA interact only with the
  helpers \texttt{warp\_atlas\_image\_into\_subject\_via\_fwd} (human)
  / \texttt{warp\_atlas\_into\_subject} (mouse, macaque) and never
  with raw forward/inverse transform lists.
\item {\em Drive calibration post-hoc.} All three species' acoustic
  forward simulations are run at unit drive $p_0=1$\,Pa and rescaled
  in post to hit a free-field-equivalent peak-negative pressure
  target (typically MI~$=1.9$ for mouse 15\,MHz, lower for human
  clinical frequencies). Relative skull transmission drops out as the
  unit-drive focal-pressure ratio in dB, independent of the calibration
  target.
\item {\em World-mm everywhere.} The slab loader and placement helper
  in all three species take apex and beam coordinates in world\,mm in
  the species' aligned-native frame, never in voxel\,mm. For the
  mouse and macaque this is the bone-centroid-centred RAS frame
  defined by PCA pre-alignment (\S5 of the mouse manuscript); for the
  human this is the NIfTI-RAS frame defined by Eq.~\ref{eq:halle_affine}
  with origin at the NRRD corner (no PCA).
\end{itemize}

\subsection{Three design choices that diverge from the mouse case}
\begin{itemize}
\item {\em Cavity-binary vs.\ intensity SyN.} The Allen CCFv3 ships a
  clean brain mask at $10$ and $25$\,$\mu$m, so the mouse pipeline
  binarises the Halle cavity at $\mathrm{BONE\_LOW}=5000$ counts and
  registers binary-to-binary with a mean-squares cost. The MNI152
  template ships only as a T1 intensity volume + a separately-derived
  brain mask; the natural Halle counterpart is the HU-windowed CT, not
  a binary cavity. Intensity-mode SyN with mutual information handles
  the cross-modality cost; the binary path is not available without
  manufacturing a Halle cavity mask that has no use elsewhere in the
  pipeline.
\item {\em Cavity centroid vs.\ outer table for placement.} The mouse
  helper picks the dorsal-most outer-bone-surface point within a 3\,mm
  cylinder around the target's (x, y) projection, then lifts it along
  $+z$ to satisfy a scalp-clearance or apex-to-target constraint.
  This is appropriate for a small mouse skull where ``dorsal'' is
  unambiguous. For the human, the head is large enough that
  ``dorsal'' is the wrong concept: a transducer placed on the side of
  the head (CP3, T3, F3) is not dorsal, but is still a valid
  transcranial placement. The perpendicular-fit placement of
  \S\ref{sec:placement} is invariant to the gross head orientation: it
  picks the surface vertex whose inward normal best aims at the target,
  not the surface vertex highest along any fixed world axis.
\item {\em No PCA pre-alignment.} The Maga 4K skull arrived in its
  scanner-axes coordinate frame, tilted $\sim$8--20$^\circ$ from RAS in
  all three Euler angles; the macaque AMNH skull similarly. The Halle
  CT, being a clinical-CT acquisition with the patient in approximate
  anatomical orientation, sits within a few degrees of RAS in
  storage. ANTs SyN's affine initialisation can absorb this residual
  rotation directly. The cost of a pre-alignment pass (order-1
  interpolation that softens cortical-bone intensities, re-anchoring of
  the acoustic ramp endpoints) outweighs the benefit; we skip it.
\end{itemize}

\subsection{One coordinate-frame migration commit}
The legacy Halle pipeline returns placement coordinates in
NRRD-voxel-mm: positions $(i\cdot v, j\cdot v, k\cdot v)$ where $i, j,
k$ are the NRRD's LAS-encoded storage indices. The TUBA convention is
Halle RAS world\,mm: positions $(-i\cdot v, +j\cdot v, +k\cdot v)$.
The migration is therefore a sign flip on the first component
(\texttt{tuba.species.human.nrrd\_voxel\_mm\_to\_halle\_ras}).
Placement positions, beam directions, surface vertices, and surface
normals all migrate by the same sign flip.

The slab basis $(\mathbf{t}_1, \mathbf{t}_2)$ does {\em not} migrate
trivially. The Gram-Schmidt seed is $(0,0,1)$ in both frames (it
points to $+S$ in both), but the cross product
$\mathbf{t}_1 = \mathrm{beam} \times \mathrm{helper}$ is a pseudo-vector
under reflections: computing the same cross product in two
right-handed coordinate systems related by a sign flip on one axis
yields {\em opposite} physical $\mathbf{t}_1$ directions. The TUBA
Halle-RAS slab is therefore the {\em physical mirror} of the legacy
NRRD-voxel-mm slab along the lateral ($\mathbf{t}_1$) axis:
$c_\mathrm{map}^\mathrm{TUBA}[i_\mathrm{lat}, j, k] =
 c_\mathrm{map}^\mathrm{legacy}[n_\mathrm{lat}-1-i_\mathrm{lat}, j, k]$.
The beam direction is the same physical direction in both frames (one
is the negative of the other in the first component, matching the
axis-sign flip); $\mathbf{t}_2 = \mathrm{beam} \times \mathbf{t}_1$
inherits the mirror from $\mathbf{t}_1$ and ends up identical in
physical direction (two pseudo-vector operations compose to a true
vector). The acoustic focal-pressure prediction is unchanged on-axis;
off-axis sidelobes are mirrored.

The regression test of the TUBA port checks both: (i) placement and
beam coordinates match the legacy values exactly after the sign-flip
conversion (max $|\Delta| = 0$ across all 11 positional/scalar fields
on the CP3$\rightarrow$S1\_left case), (ii) the slab $c\_map$ and
$\rho\_map$ match the legacy outputs to numerical tolerance after the
mirror $[\,::-1, :, :]$ along axis 0 (max $|\Delta c| = 2.4\times
10^{-4}$\,m\,s$^{-1}$ over a 0.5\,mm-pitch slab of $80\times 80\times
160$ voxels; the residual is round-off from the cross-product
re-ordering, well below visible scale).

\subsection{Lessons specific to the human migration}
Three points proved decisive in the porting work and should be carried
forward to any future clinical-CT case:

\begin{itemize}
\item {\em Trust the data, not the NRRD header.} The Halle release tags
  the volume as LPS but stores it as LAS. The disambiguation must be
  done empirically (slice-viewer inspection) and the result baked into
  a constant module (\texttt{halle\_constants.py}); downstream code
  must not re-derive the storage convention from the header.
\item {\em Outer table $\neq$ outer envelope.} The morphological
  close+fill envelope of \S\ref{sec:surface_closefill} is correct for
  registration-target generation against a head template, but it sits
  on the inner table of the bone after the 22\,mm closing inflates
  the bone mask. Placement must use the ray-cast outer table
  (\S\ref{sec:surface_raycast}), which records the outermost hit
  along each ray and is geometrically the boundary at which a
  transducer would physically sit.
\item {\em The slab basis is not canonically determined.} The
  $\mathbf{t}_1$ direction depends on the cross-product order and is
  therefore handedness-sensitive; identical physical placements
  produce mirror-flipped slabs in coordinate systems related by a
  reflection. Downstream consumers that depend on the slab basis (the
  angular-spectrum solver's per-pixel field) must either be told the
  handedness or applied to both mirror variants of the same physical
  placement and the symmetric ones declared equivalent.
\item {\em Intrinsic source-atlas asymmetry can leak past the brain
  mask after warping.} The bone-overlap diagnostic in
  \S\ref{sec:registration} caught AAL3 spilling 1.7\,\% of its label
  voxels into the cortical bone shell, 90\,\% of it on the patient-R
  side --- not a registration error, just AAL3's native R-biased
  bounding box (8\,\% more right-side voxels in its MNI source).
  Mitigation is a post-warp brain-envelope mask gated by a per-atlas
  flag (\texttt{Parcellation.mask\_to\_brain\_envelope=True}); other
  atlases stop well inside the brain mask and don't need it. The
  rule of thumb for adding a new parcellation: warp it, run the
  bone-overlap diagnostic, and turn the flag on if the overlap
  exceeds $\sim 0.5\,\%$.
\end{itemize}

\subsection{Items left for future work}
\begin{itemize}
\item {\em Live-subject CT.} The pipeline accepts any HU-calibrated
  clinical CT in NIfTI or NRRD form by changing the source path in
  \texttt{tuba.species.human.NRRD\_PATH}; nothing else needs to be
  re-tuned. The validation against a live-subject scan (with brain
  parenchyma actually present in the cavity, scalp present outside
  the skull) is straightforward to set up but has not been executed.
\item {\em Optional PCA pre-alignment.} For clinical CT acquisitions
  that arrive in non-RAS orientation (e.g., axial scans with the
  patient tilted), the mouse-style PCA pre-alignment could be enabled
  as a pre-processing pass without touching the rest of the pipeline.
  The threshold would need to be set to HU $\approx$ 700 rather than
  the mouse's uint16 30\,000, but the algorithm is unchanged.
\item {\em Bilateral target sweep.} The
  \texttt{validation\_report\_human\_sweep} dataset uses
  \texttt{ANATOMICAL\_TARGETS\_MNI} entries for 20+ targets, each
  driving an independent placement and a per-case validation report.
  A cross-case figure summarising the focal-spot metrics (lateral and
  axial FWHM, $-6$\,dB volume, intracranial PNP) across targets and
  EEG sites would close the loop with the mouse manuscript's
  Table~3-style focal-prediction comparison; this is generation-side,
  not algorithm-side, and is left to the validation-report builder.
\end{itemize}

\begin{thebibliography}{99}
\bibitem{halle_zenodo}
Halle skull microCT release, Anatomical Institute Halle, Zenodo.
Accessed 2026-04 via \texttt{setup/fetch\_halle\_skull.py}.
\bibitem{fonov2009}
V.~Fonov, A.~C.~Evans, K.~Botteron, C.~R.~Almli, R.~C.~McKinstry,
D.~L.~Collins, and the Brain Development Cooperative Group,
``Unbiased average age-appropriate atlases for pediatric studies,''
{\em NeuroImage} {\bf 54}(1), 313--327 (2011).
The ICBM 2009a non-linear symmetric template is the adult endpoint
of the same construction pipeline.
\bibitem{jurcak2007}
V.~Jurcak, D.~Tsuzuki, and I.~Dan,
``10/20, 10/10, and 10/5 systems revisited: their validity as
relative head-surface-based positioning systems,''
{\em NeuroImage} {\bf 34}(4), 1600--1611 (2007).
\bibitem{harvardoxford}
The Harvard-Oxford cortical and subcortical structural atlases,
distributed with FSL (FMRIB Software Library) and accessed via
\texttt{nilearn.datasets.fetch\_atlas\_harvard\_oxford}. Cortical
parcellation: Desikan~et~al.\ 2006; subcortical: included with
FSL by the Center for Morphometric Analysis (CMA).
\bibitem{schaefer2018}
A.~Schaefer, R.~Kong, E.~M.~Gordon, T.~O.~Laumann, X.-N.~Zuo,
A.~J.~Holmes, S.~B.~Eickhoff, and B.~T.~T.~Yeo,
``Local-global parcellation of the human cerebral cortex from
intrinsic functional connectivity MRI,''
{\em Cerebral Cortex} {\bf 28}(9), 3095--3114 (2018).
Accessed via \texttt{nilearn.datasets.fetch\_atlas\_schaefer\_2018}.
\bibitem{maga_manuscript}
G.~F.~Pinton, ``Migrating the 15\,MHz transcranial mouse pipeline
from DigiMouse to the Maga (UW) 4K microCT: orientation, cavity
extraction, and focal-prediction validation,''
\texttt{tuba/docs/mouse/manuscript.pdf} (companion report).
\end{thebibliography}

\end{document}
